\section*{Розділ IV. Прийняття рішень та голосування}
\addcontentsline{toc}{section}{Розділ IV. Прийняття рішень та голосування}

\subsection*{4.1. Процедури голосування}
\addcontentsline{toc}{subsection}{4.1. Процедури голосування}
    4.1.1. Рішення КСУ приймаються шляхом голосування делегатів КСУ відповідно до процедур, визначених цим Регламентом.

    4.1.2. Голосування на засіданнях КСУ може проводитися у таких способах:

        \begin{enumerate}[label=\alph*)]
            \item \textbf{Відкрите голосування} -- підняттям руки, поіменне голосування або інші способи, що дозволяють визначити позицію кожного делегата;
            \item \textbf{Таємне голосування} -- за допомогою бюлетенів або електронних засобів, що забезпечують конфіденційність волевиявлення;
            \item \textbf{Електронне голосування} -- з використанням спеціалізованих платформ при дистанційному проведенні засідань.
        \end{enumerate}

    4.1.3. За загальним правилом голосування проводиться відкритим способом. Таємне голосування проводиться:

        \begin{enumerate}[label=\alph*)]
            \item У випадках, прямо передбачених цим Регламентом або Положенням про ОСС;
            \item За вимогою не менше третини від присутніх на засіданні делегатів КСУ;
            \item За рішенням Спікера КСУ у разі розгляду персональних питань.
        \end{enumerate}

    4.1.4. При проведенні засідання у дистанційному форматі електронне голосування є основним способом, з обов'язковим дублюванням результатів іншими доступними засобами для забезпечення достовірності.

\subsection*{4.2. Кваліфіковані більшості}
\addcontentsline{toc}{subsection}{4.2. Кваліфіковані більшості}
    4.2.1. Рішення КСУ приймаються простою більшістю голосів від загального складу делегатів КСУ, якщо Положенням про ОСС або цим Регламентом не встановлено інше.

    4.2.2. Кваліфікованою більшістю -- не менше двох третин (2/3) голосів від загального складу делегатів КСУ -- приймаються рішення з таких питань, визначених Положенням про ОСС:

        \begin{enumerate}[label=\alph*)]
            \item Висловлення недовіри керівникам та членам органів студентського самоврядування (пункт 3.2 Положення про ОСС);
            \item Затвердження єдиного кошторису ОСС у разі відсутності погодження Студентської уповноваженої делегації (пункт 3.2 Положення про ОСС);
            \item Інші питання, для яких Положенням про ОСС встановлена вимога кваліфікованої більшості.
        \end{enumerate}

    4.2.3. У всіх випадках голосування кількість голосів "за" повинна перевищувати встановлену більшість для прийняття відповідного рішення.

\subsection*{4.3. Особливості голосування за різними категоріями питань}
\addcontentsline{toc}{subsection}{4.3. Особливості голосування за різними категоріями питань}
    4.3.1. При голосуванні з процедурних питань (затвердження порядку денного, зміни регламенту засідання, перерви тощо) достатньо простої більшості від присутніх на засіданні делегатів.

    4.3.2. При голосуванні з питань, що стосуються затвердження положень про органи студентського самоврядування, необхідна проста більшість від загального складу делегатів КСУ.

    4.3.3. При розгляді питань про висловлення недовіри посадовим особам органів студентського самоврядування голосування проводиться поіменно або таємно за вимогою особи, щодо якої розглядається питання.

    4.3.4. При затвердженні звітів органів студентського самоврядування необхідна проста більшість від загального складу делегатів КСУ.

\subsection*{4.4. Порядок підрахунку голосів}
\addcontentsline{toc}{subsection}{4.4. Порядок підрахунку голосів}
    4.4.1. При електронному голосуванні підрахунок здійснюється автоматично відповідною системою з обов'язковим дублюванням результатів секретарем засідання.

    4.4.2. При відкритому голосуванні підрахунок голосів здійснює Спікер КСУ за участю секретаря засідання.

    4.4.3. При таємному голосуванні підрахунок здійснюється лічильною комісією у складі 3 осіб, обраних КСУ з числа присутніх делегатів на початку засідання або безпосередньо перед голосуванням.

    4.4.4. Результати голосування оголошуються Спікером КСУ негайно після їх підрахунку та фіксуються у протоколі засідання з зазначенням:

        \begin{enumerate}[label=\alph*)]
            \item Загальної кількості делегатів КСУ;
            \item Кількості делегатів, присутніх на засіданні;
            \item Кількості голосів "за", "проти" та "утрималися";
            \item Результату голосування (прийнято/не прийнято).
        \end{enumerate}

\subsection*{4.5. Процедури повторного голосування}
\addcontentsline{toc}{subsection}{4.5. Процедури повторного голосування}
    4.5.1. Повторне голосування з того самого питання на одному засіданні КСУ може проводитися у таких випадках:

        \begin{enumerate}[label=\alph*)]
            \item При технічних збоях електронної системи голосування, що унеможливили участь частини делегатів;
            \item При виявленні процедурних порушень під час проведення голосування;
            \item За вимогою не менше чверті (1/4) від присутніх делегатів КСУ у разі оскарження результатів голосування з мотивованими підставами.
        \end{enumerate}

    4.5.2. Рішення про проведення повторного голосування приймає КСУ простою більшістю голосів від присутніх делегатів після заслуховування мотивів для повторного голосування.

    4.5.3. Повторне голосування проводиться не раніше ніж через 10 хвилин після оголошення результатів першого голосування для забезпечення технічної підготовки.

    4.5.4. Результати повторного голосування є остаточними та не підлягають оскарженню на тому самому засіданні.

\subsection*{4.6. Процедура при рівності голосів}
\addcontentsline{toc}{subsection}{4.6. Процедура при рівності голосів}
    4.6.1. У разі рівного розподілу голосів "за" і "проти" при першому голосуванні, Спікер КСУ оголошує перерву для додаткових консультацій тривалістю не більше 15 хвилин.

    4.6.2. Після перерви проводиться повторне голосування з того самого питання.

    4.6.3. Якщо при повторному голосуванні знову виникає рівність голосів "за" і "проти", рішення вважається неприйнятим.

    4.6.4. КСУ може прийняти рішення про перенесення розгляду питання на наступне засідання для додаткової підготовки матеріалів або консультацій.

\subsection*{4.7. Оформлення рішень КСУ}
\addcontentsline{toc}{subsection}{4.7. Оформлення рішень КСУ}
    4.7.1. Рішення КСУ оформлюється у письмовій формі та повинно містити:

        \begin{enumerate}[label=\alph*)]
            \item Назву органу, що прийняв рішення;
            \item Дату та номер засідання;
            \item Назву питання, з якого прийнято рішення;
            \item Резолютивну частину (суть прийнятого рішення);
            \item Відомості про голосування;
            \item Строки виконання рішення (за потреби);
            \item Відповідальних за виконання рішення (за потреби).
        \end{enumerate}

    4.7.2. Рішення КСУ підписується електронним підписом Спікером КСУ та секретарем засідання.

    4.7.3. Рішення КСУ набуває чинності з моменту його прийняття, якщо самим рішенням не встановлено інший строк набуття чинності.

    4.7.4. Рішення КСУ є обов'язковими для виконання всіма органами студентського самоврядування та їх посадовими особами в межах їхньої компетенції.